\documentclass[12pt]{article}
\usepackage[utf8]{inputenc}
\usepackage[T1]{fontenc}
\usepackage{hyperref}
\title{Zaghawa Language Resources for K-12 Curriculum Development}
\author{Compiled by Claude Sonnet 5.5 (claude-sonnet-5-5) for Andr\'as Kornai}
\date{}
\begin{document}
\maketitle

\section{Introduction}
Zaghawa (ISO 639-3 zag, Glottocode zagh1240), called Beria by its speakers, is a Saharan language of Sudan and Chad. Speaker estimates range from roughly 75,000 to 500,000. Several writing systems have been used.

\section{Primary Educational Resources}
No school textbooks were found in the search results. The Omniglot page mentions a Beria learning text in English and Arabic and the Beria Keyboard Project's literacy videos. I did not find these items at URLs in the search results.

\section{Linguistic and Reference Materials}
\begin{itemize}
\item Wikipedia article on the Zaghawa language: \url{https://en.wikipedia.org/wiki/Zaghawa_language}
\item Omniglot page on Zaghawa and the Beria Giray Erfe script. It mentions a Beria--English dictionary on a zaghawa4all site, which I did not verify: \url{https://www.omniglot.com/writing/zaghawa.htm}
\item Siddick Adam Issa, paper on the history of the Beria orthography: \url{https://blogs.ed.ac.uk/nilotic/wp-content/uploads/sites/3585/2021/05/Issa.pdf}
\item Resource list for non-Arabic Sudanese languages, which lists Roger Blench's Beria dictionary: \url{https://hazine.info/resources-non-arabic-sudanese-languages/}
\item Overview of the Zaghawa phonetic alphabet: \url{https://internationalphoneticalphabet.org/zaghawa}
\end{itemize}

\section{Religious and Community Resources}
No religious or community resources were found in the search results.

\section{Digital and Audio Resources}
\begin{itemize}
\item Endangered Alphabets page on the Beria script: \url{https://www.endangeredalphabets.net/?p=2750}
\end{itemize}

\section{Curriculum Development Context}
The language is used by the whole community but is not known to be taught in schools. The writing system is unsettled, with the Beria Giray Erfe script, Latin and Arabic versions all in use. A curriculum would have to choose a script first. Blench's dictionary and the Issa paper are the places to start.
\end{document}